% === VI. EXPERIMENTS AND RESULTS =============================================
\section{Experiments and Results}
\label{sec:results}

\subsection{Setup}
\label{subsec:setup}

Every run uses a \emph{static} friction surface: the effective coefficient the
identifier is scored against is \num{1.05}, constant to the last digit the
simulator reports over all \num{8186} and \num{9355} validation samples. Both
runs follow the same minimum-curvature raceline on a recorded centreline of
the CARLA silverstone circuit: a closed \SI{5828.2}{\meter} lap, reference
speed \SIrange{9.62}{25.78}{\meter\per\second}, minimum corner radius
\SI{18.7}{\meter}, peak lateral demand \SI{0.51}{\g}---roughly half the
\SI{1.0}{\g} the plant sustains. It replaces an earlier trajectory whose \SI{1.22}{\g} demand the gate of
Section~\ref{sec:integration} rejected against this plant.

Two configurations of the pipeline are run back to back on that trajectory,
three timed laps each, after an untimed out-lap to the raceline's $s=0$. They differ in four settings and nothing else: residual architecture (MLP
against S4D), objective (one-step MSE against the physics-informed loss),
Pacejka solver (trust region against differential evolution), and friction
warm start (off against on). \texttt{Baseline-NN-MSE} is the inherited configuration
evaluated at full scale on the corrected rollout, so the comparison isolates
those four settings rather than reproducing the published pipeline in full;
because they move together it is a configuration comparison and not a
per-factor ablation, and since the excitation-aware rollout is active in
\emph{both} arms it is not a test of this paper's central contribution---that
is evidenced separately in Section~\ref{subsec:sweep}.

The prior $\phip^{0}$ both arms start from, and that \eqref{eq:fitobj} pulls
toward, was obtained by fitting \eqref{eq:mf} to the simulator's per-wheel
telemetry over a dedicated limit-excitation run reaching
$\alpha=\SI{0.37}{\radian}$ and \SI{1.50}{\g}. Two consequences have
to be stated plainly: that is the same channel the force and friction results
are scored against, so the prior is not independent of the validation data; and
the manoeuvre that produced it is one the pipeline is not allowed to run, since
the raceline peaks at \SI{0.51}{\g}. On all six accepted identifications the
\texttt{Ours} arm returns the eight prior coefficients to four decimal places,
so the peak-friction accuracy below is that of a prior the fit did not move,
not of a peak factor recovered from the lap. What the corrected pipeline
demonstrably does is \emph{preserve} a physically valid peak factor where the
inherited configuration destroys it by railing $D$ on its bound.

\subsection{Identifiability of the peak factor against rollout excitation}
\label{subsec:sweep}

The claim of Section~\ref{subsec:excitation} is about the rollout configuration
alone, so it is tested where the answer is known exactly rather than in CARLA.
A tire with a known peak factor $D$ is rolled out through the synthetic sweep
of stage~2---the shipped rollout, with the residual network set to zero---and
fitted back with the shipped solver, \emph{unregularised} and started from the
true coefficients, so everything the fit loses is lost to excitation and not to
the prior, the start point or a local minimum. Rollout speed $v$, sweep amplitude $\delta_{\mathrm{sweep}}$ and true $D$ are
swept, each cell repeated over five noise realisations at the measured one-
step persistence RMSE of the two CARLA runs
($\sigma_{v_{y}}=\SI{0.009}{\meter\per\second}$,
$\sigma_{\omega}=\SI{0.005}{\radian\per\second}$). The configuration axes are
swept one at a time, since $\mu_{\mathrm{reach}}$ depends on both and the
claim that it, not the speed, governs the fit cannot be tested by moving the
speed alone. Repetition is over noise and not the solver because the noise-
free fit is seed-invariant here: five multi-start seeds and a jitter seven
times wider return the same coefficients to four decimals.

\begin{figure*}[t]
\centering
\includegraphics[width=\textwidth]{identifiability_sweep}
\caption{Peak factor $\hat D$ recovered from a known tire, unregularised and
started from the true coefficients; one series per true $D$ (dotted), mean
$\pm$ the \SI{95}{\percent} interval over five noise realisations. (a) and (b)
sweep the two arguments of \eqref{eq:mureach} one at a time, the inherited
configuration being the leftmost point of (a); starving
$\delta_{\mathrm{sweep}}$ at a fixed \SI{15}{\meter\per\second} destroys $D$
as thoroughly as starving $v$, so neither axis is a speed effect. (c) reads
both against the demand the rollout \emph{realises}, normalised by $D$: left
of the dashed line $\hat D$ reports the excitation, low where speed limits it
(filled), high where sweep amplitude does (open).}
\label{fig:sweep}
\end{figure*}

Figure~\ref{fig:sweep} says three things. First, the inherited
\SI{4}{\meter\per\second} rollout is uninformative about $D$ at \emph{every}
grip level tested: a $0.70$, a $1.05$, a $1.40$ and a $1.80$ tire all come
back as $\hat D$ between \num{0.50} and \num{0.59}, the fit reporting
$\mu_{\mathrm{reach}}$ and not the tire, with two to seven of the eight
coefficients on a bound. Second, the boundary is $\mu_{\mathrm{reach}}\gtrsim
D$ and not a speed: at \SI{6}{\meter\per\second} ($\mu_{\mathrm{reach}}=0.96$)
the $0.70$ and $1.05$ tires come back within \SI{3}{\percent} while the $1.40$
and $1.80$ tires are \SI{13}{\percent} and \SI{30}{\percent} short, and above
$\mu_{\mathrm{reach}}=1.7$ every tire is within \SI{4}{\percent} with no
systematic railing. The boundary is soft, not sharp---the $1.05$ tire is
recovered slightly \emph{below} its own $D$, a rollout stopping just short of
the peak still traversing the curvature that separates $B$, $C$ and $D$---so
the bound predicts where the information runs out. Third, the failure is not
the solver's: the noise-free fit is invariant to the multi-start seed and the
jitter width, so no solver or start point recovers a coefficient the data does
not contain.

Panel (b) also shows that the failure is not a common bias: where an under-
speeded rollout returns \emph{too little} grip, an under-swept one returns too
much---at $\delta_{\mathrm{sweep}}=\SI{0.06}{\radian}$ the \num{1.05} tire
comes back as \num{1.30} and the \num{1.80} tire rails on its bound. What the
axes share is that $\hat D$ stops being a function of the tire, the direction
of the error being set by the rollout: the signature of an unidentified
parameter rather than a biased one. They do \emph{not} collapse onto a common
$\mu_{\mathrm{reach}}$, since \eqref{eq:mureach} uses the Ackermann curvature
and tracks the realised demand only while understeer is small; against the
realised $\max|v_{x}\omega|/g$ (panel c) both behave alike.
Equation~\eqref{eq:mureach} is therefore a cheap screening bound, conclusive
when it falls below $D$, not a measure of the excitation achieved. With the
residual set to zero the structure is exactly matched, so this is a best case
and the real pipeline degrades at least this much.

\subsection{Identified tire model and peak friction}
\label{subsec:resultstire}

Figure~\ref{fig:comparison} is the cross-run comparison and
its bottom-right panel puts both identified curves against the plant's own
curve~\eqref{eq:physx}. No coefficient of \texttt{Ours} rests on a box
constraint in any of its six accepted identifications---the last accepted set
is $B,C,D,E=\num{7.076},\num{1.346},\num{1.009},\num{-2.000}$ front and
$\num{7.873},\num{1.383},\num{1.002},\num{-1.024}$ rear, inside
$B\in[4,20]$, $C\in[1.2,2.2]$, $D\in[0.4,2.0]$, $E\in[-3,1]$. Its front
$E=\num{-2.000}$ is the configured prior, not a bound, and is exactly the
coefficient the frozen Tikhonov target of \eqref{eq:fitobj} holds at its
nominal value where the lap gives no support, while the unregularised fit moves
it to $\num{-2.264}$ on no evidence.
\texttt{Baseline-NN-MSE} puts $D$ on its upper bound of \num{2.000} in the
first and the last of its six identifications and wanders over
$D_{f}\in[\num{1.241},\num{2.000}]$ in between. A peak factor of \num{2.000}
against a plant realising \num{1.05} is less a \SI{90}{\percent} parameter
error than the box-edge signature of Section~\ref{subsec:identifiability}: the
coefficient is not identified, and no per-coefficient bounds check
distinguishes that from a fit. With
the surface static, \texttt{Ours} returns $D_{f}=\num{1.009}$ and
$D_{r}=\num{1.002}$ in \emph{every} one of its six identifications, so its
$\mu$ error is a constant \num{0.041}/\num{0.048} against
\num{0.762}/\num{0.772}, \numrange{16}{19}$\times$ worse. The warm start's floor binds on no cycle here.

The price \texttt{Ours} pays is in the linear range (slopes at $\bar F_{z}=1$
and the same static axle loads throughout): \SI{13.2}{\kilo\newton\per\radian}
front and \SI{13.9}{\kilo\newton\per\radian} rear against the plant's
\num{20.4} and \num{19.0}, \SI{35}{\percent} and \SI{27}{\percent} low, where
\texttt{Baseline-NN-MSE} gives \num{19.1} and \num{18.1}, \SI{6.5}{\percent}
and \SI{4.6}{\percent} low. The understeer sign is correct in both, so both
pass the stability gate, but the baseline's margin is \SI{2}{\percent}
($l_{f}C_{f}=\SI{14.10}{\kilo\newton\meter\per\radian}$ against
$l_{r}C_{r}=\SI{14.38}{\kilo\newton\meter\per\radian}$) where \texttt{Ours}
keeps \SI{12}{\percent}: a set that lands \SI{2}{\percent} from an oversteering
pair on a coefficient it did not identify is the case the gate exists for. Only
\texttt{Baseline-NN-MSE}'s grip ceiling is a claim the plant cannot honour,
\eqref{eq:gripceiling} giving \SI{2.00}{\g} against \SI{1.01}{\g} for
\texttt{Ours} and $\approx\SI{1.0}{\g}$ measured. Both effects coexist because the lap covers only
$P_{99}(|\alpha|)=\SI{2.2}{\degree}$ against the \SI{12.0}{\degree} at which
the plant's curve peaks: it constrains the slope and says almost nothing about
the peak.

\begin{figure*}[t]
\centering
\includegraphics[width=0.41\textwidth]{tracking_error_by_scenario}\hfill
\includegraphics[width=0.41\textwidth]{mu_error_by_scenario}\\[2pt]
\includegraphics[width=0.41\textwidth]{tire_force_rmse_by_scenario}\hfill
\includegraphics[trim=0 250 0 0,clip,width=0.32\textwidth]{pacejka_identified_vs_nominal_by_scenario}
\caption{Cross-run comparison, every panel scored against the simulator's own
telemetry: lateral tracking error by active controller (the pure-pursuit boxes
precede handover and do not use the identified model); peak-factor error
against the plant's $\mu=\num{1.05}$; axle force RMSE, the one family the
inherited configuration wins, and only at the front; and the identified tire
curves against the plant's own curve~\eqref{eq:physx} at each axle's static
load (black), where \texttt{Baseline-NN-MSE}'s railed peak factor carries
roughly twice the plant's force at the edge of the range while \texttt{Ours}
tracks the peak and undercuts the slope. \textbf{Single run per configuration:
no bar or box carries a confidence interval (Section~\ref{subsec:limits}).}}
\label{fig:comparison}
\end{figure*}



\subsection{Force, prediction and closed-loop tracking}
\label{subsec:resultsforce}

The identified curve scored at the measured $(\alpha,F_{z})$ against the
plant's per-wheel force orders the arms differently across axles:
\texttt{Ours} is better at the rear (\SI{40.8}{\newton} RMSE, $R^{2}=\num{0.979}$)
and worse at the front (\SI{70.7}{\newton}, $R^{2}=\num{0.944}$, against
\SI{45.9}{\newton} and \num{0.978}), and the inherited configuration's lower
front RMSE is bought with a maximum absolute error \num{1.6}$\times$ heavier at
the front and \num{2.6}$\times$ at the rear. That
is the linear-range deficit showing up as force error, and it is the one axis
on which the inherited configuration is better: at \SI{2.2}{\degree} of slip
both curves are far below their peaks, so a railed peak factor costs nothing
while a \SI{35}{\percent} slope deficit costs \SI{71}{\newton} of front-axle
RMSE. What separates the two models is not the force at the slip angles the lap
visits but the grip ceiling those forces never probe. Neither model beats persistence at the \SI{0.035}{\second} model step ($v_{y}$
RMSE \num{0.027} against \num{0.022}, $\omega$ \num{0.032} against
\num{0.068}, persistence \num{0.009} and \num{0.005}): a one-step metric on a
smooth trajectory rewards not moving, so it is not evidence for either tire
model.

Both runs finish three timed laps with zero emergency halts and the same state
sequence---bootstrap pure pursuit, running pure pursuit, handover, running MPC
---with no fallback after handover. Under MPC, where the identified
model enters the loop, \texttt{Ours} tracks at \SI{0.092}{\meter} RMS lateral
error against \SI{0.177}{\meter}, a \SI{48}{\percent} reduction, with maximum
\SI{0.490}{\meter} against \SI{0.696}{\meter}. ITAE is neither scale-free nor comparable across runs of
different length, so two normalised forms are computed from the same samples:
$\mathrm{IAE}/T$ is \SI{0.085}{\meter} against \SI{0.132}{\meter} and
$\mathrm{ITAE}/(T^{2}/2)$ is \SI{0.059}{\meter} against \SI{0.111}{\meter}, and
all three order the runs alike. The split by controller
puts the whole difference in the MPC phase: the pure-pursuit phases give ITAE
\num{211} and \num{122}, ordered the other way and two orders of magnitude
smaller. The
QP cost is unchanged---both solve times are two orders of magnitude below the
\SI{50}{\milli\second} control period and the orderings disagree---so the
coefficients change the model the QP carries, not its size.

\subsection{Friction warm start and identification cost}
\label{subsec:mu}

The friction warm start is evaluated on a synthetic brush-tire plant rather
than in CARLA, because only a synthetic plant supplies a $\mu$ known exactly
and an excitation that can be swept independently of it. At \SI{45}{\percent}
grip utilisation---a level an ordinary lap reaches---the utilisation bound
returns \num{0.473} against a true $\mu=\num{1.05}$ while the brush fit
returns \num{1.048}, and it returns that same \num{1.048} at
\SI{63}{\percent}, \SI{80}{\percent} and \SI{100}{\percent}: flat in
excitation where the bound tracks only the grip used. Across true
$\mu\in\{0.7,1.05,1.4\}$ the fit stays within \SI{5}{\percent}, and
\SI{12}{\percent} with odometry and IMU noise; below the \SI{40}{\percent}
floor the gate suppresses it and reports the bound, so the failure is visible
rather than silent. One identification---six iterations on a \SI{30}{\second} lap---takes
\SI{20.9}{\second}, inside the re-identification period.
